\section{Discussion and Threats to Validity}
\subsection{What an improvement would establish}
Improved mapping evidence would support a bounded analysis result, not production
refactoring or quantum advantage. Reduced false confidence must be interpreted
with completion and abstention. Exact fallback can preserve behavior without
providing computational benefit.

A model experiment does not establish developer productivity; that requires a
separate user study. A bounded checker may nevertheless serve as a design debugger
without scaling to production inputs.

\subsection{Construct validity}
Keeping classical code can be justified without demonstrating mapping ability.
Separate structural, practical, and content outcomes avoid conflating these judgments.

Reference labels and candidate boundaries remain a material source of uncertainty.
Optimization, search, and validation scans can overlap in legitimate ways. A
different family or wider source region is not automatically an error. Independent
review and explicit acceptance criteria are required to interpret these choices.
The reference obligation inventory itself must not expand only for a favored
method's longer answers.

Finite tests provide bounded support. Exhausting bitstrings for selected graphs
does not exhaust all graphs or all contexts. Numeric fixtures using exactly
representable values do not prove arbitrary floating point equivalence. A local
predicate may be refuted without invalidating a complete guarded plan. Reported
outcomes must retain these boundaries, including the difference between a valid
template and a fully implemented migration.

\subsection{Internal validity}
Calls, prompts, providers, tool latency, and caches can confound comparisons.
Shared drafts, budget analyses, and time-blocked execution reduce these effects.
Unobserved service-side model changes still prevent claiming immutable weights.

The reference, interpreter, and checker may share mistakes. Independent task
definitions, equivalence controls, fault injection, and original-program comparisons
provide complementary checks. They do not prove the evaluator correct. Manual
transcription adds another error channel; passage binding and independent review
are therefore part of the planned evidence, not clerical details.

Development feedback can leak into evaluation through prompts, case selection,
checker revisions, or informal knowledge of test outputs. Separating whole
lineages and freezing the protocol limits direct leakage. Post-hoc counterexamples
remain valuable diagnostic evidence but must retain their post-hoc designation.
No repeated use of a final test is described as a fresh holdout assessment.

\subsection{External and conclusion validity}
Source-derived synthetic contexts provide control and traceability but do not
represent arbitrary industrial repositories. The limited positive families,
Python front end, small exact checks, and chosen expression language restrict
generalization. Program analysis assumptions may fail on dynamic dispatch,
reflection, complex libraries, concurrency, or native extensions. These cases
must remain unsupported rather than be counted as correctly handled.

Small or correlated source collections produce uncertain effect estimates.
Replicate calls, obligations, and test instances cannot substitute for independent
program lineages. A ceiling experiment cannot demonstrate repair effectiveness,
and absence of a measured gain does not establish that counterexamples never
help. The planned reporting preserves null results, regressions, and incomplete
delivery. Any later paper-level claim must be narrower than the evidence that
actually survives independent review.
